% page 445 of the facsimile (image p-444 of the scan)
% block 162.6 x 242.0 mm ; left margin 39.4 mm
\pgc{17.780mm}{0.000mm}{
\l{\cel{55}437.}
\l{}
\l{\sou{pia}.\cel{2}(\sou{religio}). Di qua la devoteso religiala esas vivaca,}
\l{\cel{3}forta e sincera. - EFIS.}
\l{}
\l{\sou{piafar}. (\sou{netrans.}) (\sou{aludante kavalo}). Frapar bruisoze sulo}
\l{\cel{3}per pedo, levante alte la gambi avana. - DEFS.}
\l{}
\l{\sou{piamatro}. (\sou{anat.}) Parto interna di la membrano triopla qua}
\l{\cel{3}tegas la cerebro. -}
\l{}
\l{\sou{piano}. (\sou{muziko)}. Instrumento kun klavaro, di qua la kordi ne}
\l{\cel{3}tushesas per plumo-pinto (li tale tushesis pri klavikordo),}
\l{\cel{3}ma frapesas per marteleti - procedo qua posibligas plufortigo}
\l{\cel{3}o modero di la sono segun-vole, per presar plu o min forte}
\l{\cel{3}sur la klavo qua movigas la marteleto. - DEFIRS.}
\l{}
\l{\sou{piceo}. (\sou{bot.)}\cel{3}Arboro konifera, kelke simila ad abieto. - L.}
\l{\cel{3}\sou{picea excelsa}.\cel{2}FIL.}
\l{}
\l{\sou{pichpino.(bot.}) Arboro (di Usa) di qua la bela ligno esas}
\l{\cel{3}rede-flava, riche veinoza, plu harda e plu densa kam la}
\l{\cel{3}pino-ligno ordinara, ma esas ruptebla e fendebla plu facile.}
\l{\cel{3}- L.\cel{1}\sou{pinus palustris (australis)}. - DEFIR.}
\l{}
\l{\sou{piedestalo}.\cel{2}Suportilo di kolono, di vazo, di statuo, establi-}
\l{\cel{3}sita sur sulo. - DEFIRS.}
\l{}
\l{\sou{pierido}. (\sou{zool.})\cel{2}Genero de insekti lepidoptera, ek la familio}
\l{\cel{3}\textquotedbl{}pieridei\textquotedbl{} ; un speco de oli produktas raupi qui devastas la}
\l{\cel{3}kreskaji (kauli, e c.). - L.\cel{1}\sou{pieris brassicae}. -}
\l{}
\l{\sou{pietismo}. (\sou{religio kris}t.) Nomo qua karakterizas la mento qua}
\l{\cel{3}impulsis, en la yarcento l7-esma, e mem ankore nun, grupi}
\l{\cel{3}diversa de protestantisti Germana qui tendencas ad asketismo}
\l{\cel{3}e proklamas la sacerdoteso universala di omna Deo-kredanti.}
\l{\cel{3}- DEFIRS.}
\l{}
\l{\sou{pietisto.}\cel{2}Adepto di pietismo. - DEFIRS.}
\l{}
\l{\sou{pigo.}\cel{2}(\sou{zool.)}\cel{2}Ucelo babilera, ek la familio \textquotedbl{}korvi\textquotedbl{},kun plu-}
\l{\cel{3}maro blanka e nigra, e longa kaudo dispozita etajatre. L.}
\l{\cel{3}\sou{pica}. - IS.}
\l{}
\l{\sou{pigargo}. (\sou{zool.}) Genero de raptuceli, ek la familio \textquotedbl{}falko-}
\l{\cel{3}nidei\textquotedbl{}, qua kontenas cirkum dek speci, difuzita sur la}
\l{\cel{3}Terglobo omna-regiono, ecepte en Sud-Amerika. - L.\cel{1}\sou{circus}}
\l{\cel{3}\sou{pugargus}. - EF.}
\l{}
\l{\sou{pigmeo}. I. (\sou{mitol.}) Persono di la raso \textquotedbl{}nani\textquotedbl{}. - II. Persono}
\l{\cel{3}sen merito nek valoro sociala. - DEFIRS.}
\l{}
\l{\sou{pigmento.}\cel{1}(\sou{anat.)}\cel{3}Materio ek granularo diversa-kolora, quan}
\l{\cel{3}onu trovas en la celuli di la muko-strato di epidermo, e qua}
\l{\cel{3}produktas la nuanci interdiferanta di pelo che la homo e la}
\l{\cel{3}animali. -\cel{2}DEFIRS.}
}
